\documentclass[11pt]{amsart}
\usepackage{amsmath}
\newtheorem{theorem}{Theorem}
\usepackage{hyperref}
\hypersetup{colorlinks=true,linkcolor=blue,citecolor=blue,urlcolor=blue}

\title{A Note on Badness Bounds}
\author{L. Park}
\date{}

\begin{document}
\maketitle
\tableofcontents

\begin{abstract}
We record two elementary bounds on line badness used by the incremental
page builder. Both facts are standard~\cite{knuth-plass,texbook}; the point
of this note is to fix notation.
\end{abstract}

\section{Introduction}
\label{sec:intro}

For a line whose glue must stretch by $t$ against $s$ of available stretch,
the badness is $b = 100\lvert t/s\rvert^3$, capped at $10000$. Lines with
$b \le 12$ are decent in the sense of~\cite{texbook}, and the full
paragraph optimum is characterized in~\cite{knuth-plass}. We apply the
criterion of Section~\ref{sec:bound} below on page~\pageref{sec:bound}.

\section{A uniform bound}
\label{sec:bound}

\begin{theorem}
If every line of a paragraph has badness at most $B$, then the total
demerits are at most $n(B + p)^2$ for an $n$-line paragraph with line
penalty $p$.
\end{theorem}

\begin{proof}
Each line contributes at most $(B+p)^2$ demerits by the definition quoted in
Section~\ref{sec:intro}, and demerits add across lines.
\end{proof}

The bound is sharp up to the constant, as shown by the examples
in~\cite{knuth-plass}; see also the handbook discussion~\cite{texbook}.

\begin{thebibliography}{9}
\bibitem{texbook} Donald E. Knuth. \emph{The \TeX book}. Addison-Wesley, 1984.
\bibitem{knuth-plass} Donald E. Knuth and Michael F. Plass. Breaking
  paragraphs into lines. \emph{Software: Practice and Experience},
  11(11):1119--1184, 1981.
\end{thebibliography}

\end{document}
